\section{FAIR-oriented publication and interoperability}

ECHOREPO distinguishes between its operational database structures and the
research products intended for external reuse. Internal database tables are
optimised for ingestion, validation, enrichment and application-level access,
whereas public outputs are generated as stable, machine-readable canonical
resources with documented schemas. This separation reduces the dependency of
external users on the internal implementation of the operational repository.

The publication layer exposes sample observations, images, analytical
parameters and biodiversity data as separate but related resources. These
outputs use persistent sample identifiers and include information such as
licensing, units, analytical methods and provenance-related fields where
applicable. Supporting definition tables provide additional descriptions of
parameters and analytical procedures. As a result, the internal data model can
continue to evolve without requiring downstream users to depend directly on the
operational database schema.

Canonical resources are made available through machine-oriented download
interfaces and are also packaged for long-term publication. The canonical
ECHOREPO dataset is published on Zenodo as a versioned research object
\parencite{echorepo2026dataset}. Individual releases receive persistent digital
object identifiers (DOIs), while structured metadata describe the dataset and,
where applicable, the individual tabular resources and their fields.

The publication workflow can be summarised as follows:

\begin{verbatim}
ECHOREPO operational data
          |
          v
canonical research resources
          |
          v
schema and metadata generation
          |
          v
machine-readable publication package
          |
          v
Zenodo deposition and versioning
          |
          v
persistent identifiers
          |
          v
external harvesting and discovery
          |
          v
SoilWise and other research users
\end{verbatim}

This workflow is intended to ensure that ECHOREPO data remain usable
independently of the repository's interactive web interface. A researcher or
external infrastructure can retrieve the published resources directly,
interpret their structure from accompanying metadata and documentation, and
relate records across resources using stable identifiers.

Interoperability with the wider European soil-data ecosystem is particularly
important in the context of ECHO. ECHOREPO publication packages are structured
to support discovery through SoilWise, which provides a common access and
discovery layer for soil-related research resources. For Mission Soil
resources, SoilWise recommends publication through trustworthy digital
repositories such as Zenodo and subsequently harvests the associated metadata
for discovery \parencite{soilwise2026faq}. In this model, ECHOREPO is
responsible for producing and publishing well-described research objects,
whereas broader infrastructures such as SoilWise provide aggregation and
cross-project discovery.

The implemented mechanisms can be related to the FAIR principles
\parencite{wilkinson2016fair}, as summarised in
Table~\ref{tab:fair-evidence}.

\begin{table}[htbp]
    \centering
    \caption{FAIR-oriented mechanisms implemented in the ECHOREPO publication
    workflow.}
    \label{tab:fair-evidence}
    \begin{tabular}{p{0.15\textwidth} p{0.77\textwidth}}
        \toprule
        Principle & ECHOREPO implementation and evidence \\
        \midrule

        Findable &
        Persistent identifiers for published releases, structured metadata,
        indexing in Zenodo, and discovery through external infrastructures
        such as SoilWise. \\

        Accessible &
        Machine-readable resources available through HTTPS-based download
        interfaces and persistent repository records; access conditions are
        explicitly defined for public and restricted resources. \\

        Interoperable &
        Stable tabular schemas, explicit sample identifiers, documented units
        and analytical methods, and structured metadata intended for processing
        by systems other than the ECHOREPO web application. \\

        Reusable &
        Explicit licensing, provenance-related information, documented data
        fields, versioned releases, and preservation of relationships between
        samples and derived or enriched resources. \\

        \bottomrule
    \end{tabular}
\end{table}

A key design choice is that FAIR-oriented publication is treated as part of the
normal data lifecycle rather than as a one-time export at the end of the
project. Canonical resources can be regenerated as the repository evolves,
while Zenodo versioning provides a persistent publication history. This
supports the coexistence of an actively maintained operational repository with
stable research objects that can be cited and reused independently.

Zenodo maintains both version-specific and concept-level persistent
identifiers, allowing individual releases or the evolving dataset as a whole to
be cited \parencite{zenodo2026versioning}. This distinction is useful for an
actively maintained resource such as ECHOREPO: a particular analysis can cite
the exact dataset version used, while the concept-level identifier can refer to
the dataset across successive releases.

The use of persistent identifiers and an external repository is important, but
it is not sufficient on its own to establish that a dataset is fully FAIR.
ECHOREPO is therefore described here as \emph{FAIR-oriented}. The degree to
which the published resources satisfy individual FAIR principles depends on
several factors, including metadata completeness, use of controlled
vocabularies, semantic mappings, access conditions and the ability of external
systems to interpret the published schemas.

For this reason, the evaluation presented later in the paper considers
FAIRness in terms of concrete implemented mechanisms and external
interoperability outcomes rather than treating FAIR compliance as a binary
property.